\documentclass{article}
\usepackage{amsmath,amsthm}
\newtheorem{lemma}{Lemma}
\begin{document}

\section*{Claim}
In the natural numbers, with numerals and addition defined recursively below,
\[
1+1=2.
\]

\paragraph{Strategy.} Use a direct definitional computation from the recursive definition of addition.

\section*{Setup and notation}
Let $0$ be the initial natural number and $S$ the successor operation. Define
\[
1:=S(0), \qquad 2:=S(1).
\]
Define addition recursively in its second argument by
\[
a+0:=a, \qquad a+S(b):=S(a+b)
\]
for all natural numbers $a,b$.

\begin{lemma}[Zero clause]
$1+0=1$.
\end{lemma}
\begin{proof}
Substitute $a=1$ into the defining equation $a+0:=a$. This is self-derived from the definition of addition.
\end{proof}

\section*{Main proof}
Because $1=S(0)$, the successor clause, with $a=1$ and $b=0$, gives
\[
1+1=1+S(0)=S(1+0).
\]
By Lemma 1, $S(1+0)=S(1)$. Finally, $2:=S(1)$ by definition. Therefore
\[
1+1=2.
\]
\qed

\section*{Boundary cases and audit}
The recursion reaches the base case $1+0=1$ after one successor step. All terms are natural numbers, and no hidden side conditions or target equality are used. A quantifier/definition audit and an independent circularity/base-case audit both pass.

\section*{References}
No external results are used; every inferential step is definitional or self-derived.

\end{document}
